% TOP_PROBLEM: {"problemId": "problem.infinitude-of-sophie-germain-primes", "canonicalTitle": "Infinitude of Sophie Germain Primes", "releaseRank": 390, "primaryDomain": "number_theory_arithmetic_geometry", "primaryDomainLabel": "Number theory & arithmetic geometry", "displayStatus": "open", "releaseStatus": "open"}
% !TeX program = lualatex
% Definition and sources only; this file is not an LLM solution attempt.
\documentclass[11pt]{article}
\usepackage[margin=25mm]{geometry}
\usepackage{fontspec}
\setmainfont{Latin Modern Roman}
\usepackage{amsmath,amssymb,mathtools}
\usepackage{unicode-math}
\setmathfont{Latin Modern Math}
\usepackage{xurl,hyperref}
\hypersetup{colorlinks=true,urlcolor=blue}
\setcounter{secnumdepth}{0}
\setlength{\parindent}{0pt}
\setlength{\parskip}{5pt}
\setlength{\emergencystretch}{3em}
\begin{document}
\section*{\texorpdfstring{Infinitude of Sophie Germain Primes}{Infinitude of Sophie Germain Primes}}\label{390}
\subsection{Definitions and mathematical statement}
A Sophie Germain prime is a positive prime $p$ for which $2p+1$ is also prime. The infinitude question is
\[
 \forall X>0\ \exists p>X:\quad p\text{ and }2p+1\text{ are prime}.
\]
Equivalently, the number of such $p\le x$ should tend to infinity as $x\to\infty$. No asymptotic counting formula, density constant, or effective bound for the next example is required. Primality of both linear expressions is essential; replacing one with an integer having two prime factors is only a weaker approximation to the question.
\par\smallskip\textit{Formulation and concept references:} \href{https://mathworld.wolfram.com/SophieGermainPrime.html}{[S1]}.

\subsection{Short English statement}
Are there arbitrarily large primes $p$ for which $2p+1$ is also prime?
\subsection{Sources}
\begin{itemize}
\item[S1]Sophie Germain Prime (MathWorld).
\url{https://mathworld.wolfram.com/SophieGermainPrime.html}.
\textit{Formulation source.}
\end{itemize}
\end{document}
